\Needspace{8\baselineskip}
\section{Harness与业务应用：同事最终使用什么}

\subsection{从模型接口到员工工作台}
非编码工作同样需要模型之外的能力。读取财报涉及文件解析与检索，分析表格涉及计算，整理研究材料涉及引用、文件生成和任务记录。Harness将模型与这些工具组织起来，完整应用再提供同事能够直接使用的界面。

因此，应用层可以按最终工作形态理解：员工工作台提供日常入口，资料研究平台连接知识与证据，流程平台处理重复工作，通用助手围绕目标持续行动，开发框架则用于构建更贴合内部系统的应用。

\begin{table}[H]\small\centering
\caption{AI应用的主要形态与业务价值}\label{tab:application-types}
\begin{tabularx}{\linewidth}{@{}P{28mm}Y Y@{}}\toprule
应用形态 & 代表项目 & 使用者得到的能力\\\midrule
员工AI工作台 & Open WebUI、LibreChat、AnythingLLM & 在统一入口聊天、上传文件、使用资料与专用助手。\\
资料研究平台 & Onyx、RAGFlow & 在大量资料中查找证据、比较信息、形成可追溯回答。\\
固定流程平台 & Dify、n8n、Flowise & 把一次有效操作固化为可以重复运行的处理流程。\\
通用任务助手 & Goose、OpenClaw & 接收一个目标，组合文件和工具操作，持续交付结果。\\
应用开发框架 & LangGraph、Microsoft Agent Framework、Agno、CrewAI & 由工程团队定义业务步骤、工具、状态和人工处理。\\\bottomrule
\end{tabularx}
\source{根据官方产品定位与功能归纳，类别可重叠。\src{TECH-043}\src{TECH-044}\src{TECH-045}\src{TECH-030}\src{TECH-047}\src{TECH-049}}
\end{table}

\subsection{贯穿示例：两期财报与一份CSV形成比较表}\label{sec:case}
以下使用虚构公司和数据说明工作过程。输入包括云川公司2024年、2025年财报各一份，以及按年、季度记录营业收入的CSV。两份财报均以人民币亿元列示持续经营业务，营业收入分别为120、138，营业成本分别为84、96.6；CSV中两年的季度收入分别为28、30、29、33和32、34、35、37。

任务是形成带期间、单位和来源的比较表，并给出三条说明。模型识别财报字段与注释，计算工具完成单位对齐、加总和比率计算，再由应用组织引用与结果文件。表\ref{tab:case-output}展示所需产物。

\begin{table}[H]\small\centering
\caption{财报比较任务的说明性输出}\label{tab:case-output}
\begin{tabularx}{\linewidth}{@{}P{30mm}P{20mm}P{20mm}Y@{}}\toprule
指标 & 2024年 & 2025年 & 来源及计算依据\\\midrule
营业收入（亿元） & 120 & 138 & 各期财报第12、14页；CSV季度加总分别一致。\\
营业成本（亿元） & 84 & 96.6 & 各期财报同页，币种、单位及业务范围一致。\\
毛利率 & 30.0\% & 30.0\% & （营业收入减营业成本）除以营业收入。\\\bottomrule
\end{tabularx}
\smallnote{示例设定：公司、数字和页码均为虚构；金额统一为人民币亿元。}
\end{table}

配套说明为：营业收入增长15\%；营业成本同比例增长，毛利率保持30\%；两年CSV的季度加总分别为120和138亿元，与财报一致。这个任务把材料理解、确定性计算和来源核对组成一个完整产物。下文几种产品形态的差别，主要体现在谁来安排这些步骤。

\subsection{员工工作台：聊天、文件与可复用助手}
\subsubsection{Open WebUI}
Open WebUI提供自托管的多模型工作台。员工通过浏览器选择模型、上传资料、使用知识库或调用已经配置的工具，后台连接则由统一平台管理。其官方功能包括用户组、角色权限、模型配置和可扩展工具。\src{TECH-043}

这一形态的优势在于把分散的模型使用方式集中到一个入口。研究员可以围绕公开财报连续提问，运营人员可以查询允许访问的制度资料，团队还可以保存面向不同工作的专用配置。模型服务与使用界面分离，也便于后台更换模型。

\subsubsection{LibreChat与AnythingLLM}
LibreChat将多模型聊天、专用Agent、文件搜索和工具操作结合。用户可以为一个助手保存指令、资料和工具，并将其共享给同事。其当前文档还覆盖文件上传、代码执行与生成结果下载。\src{TECH-044}

在上述比较任务中，使用者可以先上传三份文件，要求列出指标，再追问增长率、要求核对CSV，最后下载结果；任务方向由连续提问逐步确定。LibreChat可将文件搜索、计算与下载接在同一会话中，其优势是保留探索和修正空间。代码执行通过可自托管的计算服务提供；持久会话和附接工作环境则是当前标为实验性的新功能。\src{TECH-044}（Agents、Code Interpreter文档）

AnythingLLM将模型连接、资料导入、文档聊天与Agent能力整合在一起，提供桌面和自托管方式。围绕行业或项目建立资料工作空间后，使用者可以持续整理材料、查询信息和形成草稿。其优势是个人资料与模型入口比较集中；当前README将多人及权限能力对应到Docker路线，桌面、团队部署和Pro产品提供的管理功能不同。\src{TECH-053}

\subsection{资料研究平台：让回答能够回到证据}
\subsubsection{Onyx：连接企业资料的搜索与研究助手}
Onyx连接已有信息源，提供搜索、问答、专用Agent、多步研究和文件产物。它既包含面向员工的使用入口，也包含完整部署所需的资料索引与检索组件。\src{TECH-045}

对于基金研究，这种产品形态可以围绕公司或行业持续查找依据。例如，系统取得允许访问的材料，比较不同时间的观点，再组织为带证据的研究简报。其主要价值在于将资料发现与研究问题放在同一工作过程，减少逐个平台寻找信息的时间。

企业资料的访问范围会随人员和文档变化，因此连接器和权限继承是这类系统的重要组成。Onyx提供社区核心和企业功能，企业身份与权限等能力按相应产品版本提供，企业代码目录采用单独许可。\src{TECH-045}（LICENSE及产品版本页面）

\subsubsection{RAGFlow：解析复杂文档并组织可追溯回答}
RAGFlow重点处理资料解析、切分、索引和基于证据的回答，并将检索与Agent流程结合。对于长PDF、跨页表格或包含注释的财报，资料如何被转换为模型上下文，会直接影响最终答案。\src{TECH-046}

RAG流程可以概括为：系统先将文档整理成可检索资料；使用者提问后，检索组件选取有关片段；模型据此形成回答，应用保留来源位置供核查。RAGFlow的特色在于资料处理与引用检查，而Open WebUI、Dify等产品也可以利用相同机制提供知识问答。

\begin{figure}[H]\centering
\flowfive{原始资料\\财报、公告、制度}{文档处理\\解析、切分、索引}{查找证据\\问题与相关片段}{模型回答\\归纳、比较、解释}{结果核查\\引用与原文位置}
\caption{资料问答的共同流程。不同产品的差异主要体现在文档处理、资料连接、检索与用户核查方式。}
\end{figure}

Onyx与RAGFlow的侧重点由此形成区别：前者强调连接企业知识并提供研究入口，后者突出文档成为高质量上下文的过程。在前述例子中，如果财报尚未收集，Onyx类平台侧重从已有资料源找到正确期间的报告；当PDF已经取得，RAGFlow类平台侧重保留跨页表格、单位与脚注，便于回答回到原文。本例的三份指定材料可直接解析后交给模型；资料持续积累时，索引和检索使已有材料能够重复利用。

\subsection{固定流程平台：把重复工作变成应用}
\subsubsection{Dify：通过可视化流程提供部门应用}
Dify将模型、资料、工具和判断条件放在可视化工作流中，可以形成聊天助手或固定业务应用。其价值是让处理过程可见，并将模型配置、知识库、应用入口及运行记录集中起来，便于业务人员与工程团队讨论同一套流程。\src{TECH-030}

同一财报比较任务可以预先固定为：接收两份报告和CSV，提取年度、币种、单位、收入与成本，程序计算增长率和毛利率，并核对季度加总。若期间或单位无法对齐，流程将疑点和出处交给人员确认，再保存比较表。Dify提供搭建与运行这些步骤的平台，应用配置明确字段、核验规则、导出和复核入口，使使用者可以重复运行同一方法。

\begin{figure}[H]\centering
\flowfive{接收材料\\两期财报与CSV}{模型提取\\期间、单位与字段}{程序计算\\加总、增长与毛利率}{人员确认\\疑点与出处}{保存产物\\比较表与三条说明}
\caption{财报比较的固定流程设计：从接收材料、提取和计算，到人员确认及结果保存。}
\end{figure}

Dify自托管还包含数据库、缓存、插件、后台任务进程及隔离的计算环境。可视化搭建整合界面与流程，后台组件负责状态保存、任务执行和计算隔离；多工作空间和品牌条件见附录\ref{app:software-license}。\src{TECH-030}

\subsubsection{n8n与Flowise：业务连接和可视化搭建}
n8n侧重连接现有业务系统，通过触发器把输入、规则、模型处理和后续动作串起来。例如，新资料进入后可以依次完成归类、整理和归档，模型参与其中需要语言理解的环节。它的优势主要体现在跨系统自动化及已有连接器。\src{TECH-034}

Flowise通过Chatflow与Agentflow展示模型、资料和工具之间的关系，便于试验不同处理流程。与Dify相似，它为搭建应用提供可视化入口，团队管理和企业功能按版本与许可提供，汇总见附录\ref{app:software-license}。\src{TECH-033}\src{TECH-034}

这些平台改变的是工作组织方式。它们把重复的人工操作转成可运行、可检查的过程，模型负责语言理解，资料、规则和工具提供处理依据与执行能力。

\subsection{通用任务助手：围绕一个目标持续完成工作}
\subsubsection{Goose与OpenClaw}
Goose提供桌面应用、CLI和API，通过模型与工具完成研究、写作、自动化和数据分析。其当前官方项目为aaif-goose/goose，具有自定义模型服务和扩展能力，并以Recipes保存可复用的指令、工具及设置。\src{TECH-047}

对于同一比较任务，研究员可以一次提出“读取这两期财报和CSV，生成比较表及三条说明”。Goose依据目标选择读取、计算和写文件的步骤；这些能力通过配置的文件、解析和执行工具提供。处理顺序由助手在任务中决定，便于临时增加分析问题或更换产物格式。桌面入口提供交互界面，Recipes保存有效的方法，扩展机制连接所需的文件和业务工具。

OpenClaw更强调持续个人任务。它将会话、记忆、工具和不同入口连接到自托管Gateway，使关注主题和工作背景可以延续到多次交互。资料整理、会议准备和持续跟踪因此可以围绕同一个助手开展。\src{TECH-048}

两者的使用形态各有重点：Goose提供面向电脑任务的通用操作入口，OpenClaw强调持续状态与多渠道访问。OpenClaw官方以一个Gateway对应一个信任边界；当不同人员需要严格隔离时，其个人助手配置与企业共享平台存在不同的管理要求。\src{TECH-048}（Trust model）

\subsection{开发框架：把业务逻辑落实为专用系统}
现成应用提供使用入口，开发框架则让团队进一步掌握处理步骤、工具和任务状态。对于具有特定数据接口、审批过程或输出格式的工作，框架决定了这些规则怎样进入实际应用。

\begin{table}[H]\small\centering
\caption{开发框架的主要作用与优点}
\begin{tabularx}{\linewidth}{@{}P{35mm}Y Y@{}}\toprule
项目 & 做什么 & 主要价值\\\midrule
LangGraph & 组织步骤、分支、任务状态与人工中断。 & 过程清楚，可把固定流程与开放研究结合。\\
Microsoft Agent Framework & 提供Agent与多步流程的统一开发基础。 & 连接多模型、工具与企业应用的路线较完整。\\
Agno & 组合Agent、协作团队、流程及运行管理。 & 将知识、记忆和助手服务组织在同一产品路线。\\
CrewAI & 按职责组织Agent，并用Flow控制过程。 & 资料搜集、核对和写作等分工较直观。\\
Haystack & 将检索、排序、过滤和生成组成处理流程。 & 资料处理和答案依据较容易分环节检查。\\
LlamaIndex / LangChain & 提供文档、索引、模型及工具连接组件。 & 复用现有集成，减少重复建设。\\\bottomrule
\end{tabularx}
\source{官方README与概念文档。\src{TECH-029}\src{TECH-049}\src{TECH-051}\src{TECH-052}\src{TECH-032}\src{TECH-031}\src{TECH-028}}
\end{table}

LangGraph组织状态、分支和人工中断，LangChain提供模型与工具连接。Deep Agents进一步组合规划、文件工具、上下文管理和子任务，让工程团队复用多步助手的常见能力。这些能力为开发专用研究助手提供了现成组件。\src{TECH-028}\src{TECH-029}\src{TECH-041}

DeepSeek Harness同样提供可组合的模型、工具、会话和执行循环，适合定制任务组织方式。当前项目处于开发者预览阶段，重点提供可组合的执行机制，工程团队在其上建设业务入口、数据连接和运行管理。\src{TECH-037}

Microsoft Agent Framework提供.NET和Python等开发路线，包含工具、工作流、人工处理和多模型服务连接。微软官方说明其1.0已于2026年4月正式发布；AutoGen当前进入维护模式，并将新项目引向Agent Framework。本地模型连接使其可以用于内部服务，云端托管产品则提供另一种运行方式。\src{TECH-049}\src{TECH-050}

Agno当前区分SDK、AgentOS运行时与Control Plane管理界面，适合把多个专用助手作为持续服务管理。CrewAI则通过Crews表达协作分工，通过Flows控制整体处理顺序。多Agent适合将资料搜集、数字核对和写作拆成独立工作，再汇总为完整产物；模型调用和结果整合构成相应的运行成本。\src{TECH-051}\src{TECH-052}

Haystack和LlamaIndex更接近资料与检索工作。其中Haystack强调显式组件与流程，LlamaIndex当前公开方向更加重视文档解析及相关产品。开放组件提供可复用的开发能力，商业管理平台和托管服务则承担额外的运行工作。\src{TECH-032}\src{TECH-031}

\subsection{编码Agent与研究人员的文件工作}
OpenCode、pi和Codex CLI提供模型与工具结合的编码工作环境，也可以用于技术研究人员的资料、脚本和文件任务。OpenCode侧重较完整的交互体验，pi侧重精简与可扩展性，Codex CLI提供开源的本地Agent与自定义模型服务配置。\src{TECH-035}\src{TECH-036}\src{TECH-038}

这类工具把脚本、文件和计算环境放在同一工作空间，适合技术研究人员直接处理数据与产物。面向部门同事的工作台则进一步提供浏览器界面、现成文件管理和协作入口。

Claude Code则属于商业产品。当前公开仓库的许可保留所有权利，Anthropic官方也明确不支持通过gateway将其路由到非Claude模型。SGLang等引擎已有相关技术适配文档；软件许可与官方支持范围见附录。\src{TECH-039}\src{TECH-021}
